% =====================================================================
% SHORT RESEARCH PROSPECTUS — Where Does the Solenoid Burden Go?
% Build: cd research-ideas/centre-stack && latexmk -pdf research-prospectus.tex
% =====================================================================
\documentclass[11pt]{article}
\usepackage[letterpaper,margin=0.85in,top=0.75in,bottom=0.9in]{geometry}
\usepackage[T1]{fontenc}
\usepackage[scaled=0.95]{helvet}
\renewcommand{\familydefault}{\sfdefault}
\usepackage{amsmath}
\usepackage{xcolor}
\usepackage[most]{tcolorbox}
\usepackage{enumitem}
\usepackage{fancyhdr}
\usepackage[colorlinks=true,linkcolor=navy,urlcolor=navy]{hyperref}

\hypersetup{
  pdftitle={Where Does the Solenoid Burden Go?},
  pdfauthor={Jony (TUNEM)},
  pdfsubject={Six-week feasibility study of merging-compression startup and NI-HTS centre-stack compatibility},
  pdfkeywords={spherical tokamak, NI-HTS, centre stack, startup, research prospectus}}

% ---- house palette (shared with the illustrated guide and excerpt)
\definecolor{navy}{HTML}{1B3A5C}
\definecolor{gold}{HTML}{D4A028}
\definecolor{teal}{HTML}{1F7A8C}
\definecolor{cream}{HTML}{FFF8E7}
\definecolor{mist}{HTML}{EEF3F7}
\definecolor{ink}{HTML}{222222}
\definecolor{mute}{HTML}{666666}

\setlength{\parindent}{0pt}
\setlength{\parskip}{6pt}
\color{ink}

% ---- running header (pages 2–3) and footer
\pagestyle{fancy}
\fancyhf{}
\renewcommand{\headrulewidth}{0pt}
\fancyfoot[L]{\footnotesize\color{mute}Research prospectus \,|\, discussion draft}
\fancyfoot[R]{\footnotesize\color{mute}\thepage}
\newcommand{\runninghead}[1]{%
  \begin{tikzpicture}[remember picture, overlay]
    \fill[navy] (current page.north west) rectangle ([yshift=-0.62in]current page.north east);
    \node[anchor=west, text=white, font=\footnotesize\bfseries] at ([xshift=0.85in, yshift=-0.33in]current page.north west) {WHERE DOES THE SOLENOID BURDEN GO?};
    \node[anchor=east, text=white, font=\footnotesize] at ([xshift=-0.85in, yshift=-0.33in]current page.north east) {#1};
  \end{tikzpicture}\par\vspace*{0.45in}}

% ---- building blocks
\newcommand{\kicker}[1]{\noindent{\footnotesize\bfseries\color{gold}\MakeUppercase{#1}}\par\vspace{4pt}}
\newcommand{\kickerteal}[1]{\noindent{\footnotesize\bfseries\color{teal}\MakeUppercase{#1}}\par\vspace{4pt}}
\newcommand{\bigtitle}[1]{{\fontsize{26}{30}\selectfont\bfseries\color{white}#1\par}}
\newcommand{\subtitle}[1]{{\normalsize\color{white}#1}\par}
\newcommand{\pagetitle}[1]{{\fontsize{20}{24}\selectfont\bfseries\color{navy}#1\par}\vspace{6pt}}
\newcommand{\heading}[1]{\vspace{8pt}{\large\bfseries\color{navy}#1}\par\vspace{2pt}}
\newcommand{\subheading}[1]{\vspace{6pt}{\normalsize\bfseries\color{navy}#1}\par\vspace{1pt}}

\newtcolorbox{goldbox}[1]{enhanced, colback=cream, colframe=gold, boxrule=0.8pt,
  arc=3pt, left=8pt, right=8pt, top=6pt, bottom=6pt, fonttitle=\bfseries,
  coltitle=navy, title=#1, attach title to upper={\par\vspace{2pt}}, before skip=8pt, after skip=8pt}
\newtcolorbox{mistbox}[1]{enhanced, colback=mist, colframe=mist, boxrule=0pt,
  arc=3pt, left=8pt, right=8pt, top=6pt, bottom=6pt, fonttitle=\bfseries,
  coltitle=navy, title=#1, attach title to upper={\par\vspace{2pt}}, before skip=8pt, after skip=8pt}

\setlist[itemize]{label=\textcolor{teal}{\textbullet}, leftmargin=1.4em, itemsep=2pt, topsep=3pt}
\setlist[description]{font=\bfseries\color{teal}, labelwidth=4.2em, leftmargin=5.2em, itemsep=2pt, topsep=3pt, style=sameline}

\begin{document}

% ======================================================= PAGE 1
\begin{tikzpicture}[remember picture, overlay]
  \fill[navy] (current page.north west) rectangle ([yshift=-2.05in]current page.north east);
  \fill[gold] ([yshift=-2.05in]current page.north west) rectangle ([yshift=-2.12in]current page.north east);
\end{tikzpicture}%
\vspace*{-0.15in}
\kicker{Short research prospectus \,|\, September 2026}
\bigtitle{Where Does the Solenoid Burden Go?}
\vspace{4pt}
\subtitle{A six-week feasibility study of merging-compression startup and NI-HTS centre-stack compatibility}
\vspace{1.15in}

\kickerteal{Primary question}
{\Large\bfseries\color{navy}How much central-solenoid burden can merging-compression genuinely remove, and where does that burden reappear?\par}
\vspace{4pt}

The proposed first project connects two existing research streams: (1) sharp-D no-insulation
high-temperature-superconducting toroidal-field windings, and (2) merging-compression and repetitive startup
in spherical tokamaks. The initial task is not a complete reactor optimization. It is a bounded test of whether
one representative startup transient creates a consequential magnet and recovery penalty.

\begin{goldbox}{Burden-transfer hypothesis}
Reducing solenoid flux does not automatically remove startup cost. The burden may reappear as induced current,
HTS hysteresis, structural eddy loss, TF field error, cryogenic recovery time, actuator power, failed formation
shots, cyclic damage, or reduced plant duty.
\end{goldbox}

\heading{Proposed first deliverable}
\begin{itemize}
  \item Reproduce one published charging/discharging response of the sharp-D NI-HTS winding.
  \item Apply one representative or sanitized CS/PF/merging startup waveform.
  \item Calculate the field seen by the winding, the field error, and the loss-channel decomposition.
  \item Issue a one-page verdict: signal, obstruction, indeterminate, or irrelevant.
\end{itemize}

\begin{mistbox}{Scope discipline}
One primary proposal, with later branches activated only by the first result. No claim of novelty is made until
the literature and group-internal work are checked.
\end{mistbox}

\clearpage

% ======================================================= PAGE 2
\runninghead{Method and kill tests}
\kickerteal{A bounded first calculation}
\pagetitle{Six-week feasibility sprint}

\begin{description}
  \item[Week 1] Reproduce a published NI-HTS charging/discharging case and document model boundaries.
  \item[Week 2] Define one startup cycle and compute the time-dependent field at the winding.
  \item[Week 3] Add minimal vessel/casing/passive-current paths and check energy conservation.
  \item[Week 4] Decompose cold energy by HTS hysteresis, structure/casing eddy, contact, coupling, and joints.
  \item[Week 5] Sweep only the parameters that move the dominant channel and the useful plasma-field transfer.
  \item[Week 6] Report the verdict, uncertainty, missing data, and whether a paper-scale model is earned.
\end{description}

\heading{Quantities to report}
\begin{itemize}
  \item Useful startup field at the plasma and unwanted TF-field perturbation.
  \item Cold joules per event, local temperature-rise proxy, and recovery time.
  \item Dominant loss channel and its sensitivity to waveform, contact, geometry, and passive structures.
  \item The central-solenoid flux or radial space actually saved at the same post-startup plasma target.
\end{itemize}

\begin{goldbox}{Kill tests and honest pivots}
\begin{itemize}
  \item Forcing at the winding is negligible: stop the magnet-coupling paper.
  \item HTS hysteresis dominates: study tape orientation, field angle, actuator placement, or slew.
  \item Casing/structure eddy loss dominates: study segmentation and field penetration.
  \item Contact-controlled response is material: test a small number of manufacturable resistance zones.
  \item Plasma startup uncertainty dominates: pivot to formation reliability and equivalent flux saved.
\end{itemize}
\end{goldbox}

\clearpage

% ======================================================= PAGE 3
\runninghead{Collaboration and paper path}
\kickerteal{Conditional extensions}
\pagetitle{What the first result could unlock}

\subheading{If the coupling is consequential}
Compare CS-heavy, hybrid, and merge-heavy architectures at one fixed plasma mission. Parameterize startup
allocation by $x_{\mathrm{CS}} = \Phi_{\mathrm{CS}}/\Phi_{\mathrm{required}}$ and map feasibility in
$(x_{\mathrm{CS}},\ \text{repetition rate})$. Report millimetres recovered, cold energy, field error, recovery,
duty, and the active constraint.

\subheading{If selective transmission appears possible}
Test whether actuator placement, casing segmentation, tape orientation, winding topology/contact architecture,
and waveform co-design can increase useful plasma field per unit cold-mass loss and TF-field distortion.

\subheading{If a robust obstruction appears}
Ask whether passivity or reciprocity imposes a quantitative limit on transmitting the startup band while
suppressing loss and field error. This theory branch activates only after the numerical model displays the
obstruction.

\begin{mistbox}{Inputs requested from Professor Tan}
\begin{itemize}
  \item The intended operating cycle: what repeats, what remains energized, and what sets dwell.
  \item One representative or normalized startup waveform and approximate actuator geometry.
  \item The metric that actually limits the NI-HTS winding or plasma scenario.
  \item Whether external mutual forcing can be added to the existing equivalent-circuit model.
  \item What work is already complete internally, and the publication/IP boundary.
\end{itemize}
\end{mistbox}

\heading{Evidence base to reproduce first}
{\small
Yin, Tan et al., \emph{Fusion Engineering and Design} 220, 115338 (2025);
Yin et al., \emph{IEEE Transactions on Applied Superconductivity} 36, 4200305 (2026);
Tan, Gao et al., IAEA FEC 2025 (SUNIST-2);
Yang et al., \emph{Superconductor Science and Technology} 38, 055008 (2025).\par}

\end{document}
